\section{方法}
\label{sec:method}

この章は「何をしたか」と同時に「なぜそうしたか」を述べる。
特に \ref{sec:circular} 節から \ref{sec:neff} 節までは、
結果の図の読み方に直結するので飛ばさないでほしい。

\subsection{時間の切り方}
\label{sec:weeks}

1 か月を次の 2 つに分ける。

\begin{center}
\begin{tabular}{lll}
\toprule
ビン & 日 & 日数 \\
\midrule
H1 & 1--15 日 & 15 \\
H2 & 16 日--月末 & 13--16 \\
\bottomrule
\end{tabular}
\end{center}

1 年で 24 ビンになる。

\begin{whybox}{なぜ ISO 週や半旬ではなく「月内固定」なのか}
境界の決め方には候補が 3 つあった。

\textbf{ISO 週}（月曜始まりの通し番号）は年ごとに月との対応がずれる。
「1 月の 2 番目のビン」と言ったとき、年によって指す日付が変わってしまう。

\textbf{半旬}（5 日 $\times$ 73）は気候学では標準的で、標本数が完全に揃う。
しかし月境界をまたぐので、「何月の話か」を言うのに一手間かかる。

\textbf{月内固定}は、「3 月前半」が毎年必ず 3 月 1--15 日を指す。
打ち上げ計画は「来年の 3 月前半に上げたい」という形で立つものなので、
人が計画を立てる単位にそのまま対応する。これを選んだ。
\end{whybox}

\begin{whybox}{なぜ 7 日ではなく 15 日なのか}
本報告は当初、1 か月を 7 日 $\times$ 4 に分けていた（W1 = 1--7 日、
W2 = 8--14 日、W3 = 15--21 日、W4 = 22 日--月末）。
7 日という幅は、データからどこまで細かく区別できるかを調べる前に決めたものである。

第 \ref{sec:validation} 章がそれを実測した。1 年抜き交差検証の誤差は
\numrange{5}{15} 日のあいだで平坦で、1 標準誤差ルールが選ぶのは和歌山と愛媛で
\textbf{15 日}、北海道で 30 日である。
7 日は間違いではない。「区別できない範囲」の中に入っている。
ただし\textbf{データから読み取れる細かさよりも細かかった}。

\textbf{精度は変わらない。} 季節ごとに交差検証をやり直すと、
15 日ビンと 7 日ビンの差は北海道・和歌山の 2 領域 $\times$ 4 季節の 8 通りすべてで
1 標準誤差未満に収まり、うち 6 通りでは 15 日のほうがわずかに良い。

\textbf{変わるのは分位点の精度である。} 6 時間ごとの解析は互いによく似ていて
（ラグ 1 自己相関がおよそ 0.975）、7 日ビンの 280 標本は実質 4 個ぶんの
情報しか持たない（\ref{sec:neff} 節）。
ビンを倍にすると、本報告が主に使う量のノイズがほぼ半分になる。
\SI{200}{\hecto\pascal} の風速 p90 の標準誤差は、
7 日で \SI{1.70}{\metre\per\second}、15 日で \SI{1.04}{\metre\per\second} である。

7 日の製品は \ref{sec:appendix-week} に残してある。
\end{whybox}

代償は H2 だけが 13--16 日と変動することである。
H1 は常に 15 日なので標本数は 10 年で 600 個、
H2 は月によって 520--640 個になる。
したがって\textbf{ビンごとに標本数が違う}。
これを忘れると、たとえば「H2 だけ最大風速が大きい」という見かけの結果が出る
（標本が多いぶん極値を引きやすいだけ）。
このため標本数 $n$ はすべての統計量と並べて保存し、
各領域の章の最初に示す（図\ref{fig:f0-hokkaido}、\ref{fig:f0-wakayama}）。

なお 7 日ビンの日数は 7--10 日で 43 \% 変動していたので、
この点は以前より改善している。

\subsection{気圧面を高度に直す}
\label{sec:altitude}

データは気圧面上にあるが、人が読みたいのは高度である。
換算には 2 段階あり、混同しやすい。

\paragraph{第 1 段階：ジオポテンシャル高度}
JRA-3Q が持っている $Z$ は\textbf{ジオポテンシャル高度}で、単位は gpm（ジオポテンシャルメートル）である。
これは距離ではなく、重力に逆らって単位質量を持ち上げるのに要した仕事を
標準重力 $g_0 = \SI{9.80665}{\metre\per\second\squared}$ で割った量である。

\begin{equation}
  Z = \frac{1}{g_0}\int_0^z g(\phi, z')\, \mathrm{d}z'
\end{equation}

\paragraph{第 2 段階：幾何高度}
実際の重力は高度とともに弱まり、緯度によっても違う。
そのためジオポテンシャル高度と実際の距離（幾何高度 $z$）はずれる。

\begin{equation}
  z = \frac{R_\mathrm{eff}(\phi)\, Z\, g_0}{g(\phi)\, R_\mathrm{eff}(\phi) - Z\, g_0}
\end{equation}

ここで $g(\phi)$ は Somigliana の式による海面重力、
$R_\mathrm{eff}(\phi)$ は List の有効地球半径 \citep{list1968smithsonian} である。

補正量は \ang{45}N・\SI{30}{\kilo\metre} で約 \SI{+140}{\metre}。
一見わずかだが、上層の気圧面どうしの間隔と同じ桁であり、無視すると層が 1 枚ずれかねない。

\begin{whybox}{なぜ変数名を \texttt{height} にしないのか}
コード中では必ず \texttt{\_gpm}（ジオポテンシャル）か \texttt{\_geom\_m}（幾何）を付けている。
単に \texttt{height} と書くと、どちらの意味で使われているか 3 か月後には分からない。
この取り違えは古典的な事故であり、名前で防ぐのが一番安い。
\end{whybox}

なお既存のシミュレータ（\texttt{backend/bfps/core/atmosphere.py}）は、
気圧から高度への換算に国際標準大気 ISA \citep{icao1993isa} を使っている。
ISA は季節も緯度も天気も知らない理想化なので、実際とどれだけ違うかは確かめる価値がある。
そのずれは図\ref{fig:f10-hokkaido} と図\ref{fig:f10-wakayama} に示す。

\subsection{なぜ風向を平均してはいけないのか}
\label{sec:circular}

風向は角度であり、\ang{0} と \ang{360} は同じ向きを指す。
つまり数直線ではなく円の上にある。
円の上の量に、ふつうの算術平均を当てると誤った答えが出る。

\begin{whybox}{北風 2 回の平均が南風になる}
ほぼ真北から吹く風を 2 回観測したとする。1 回目は \ang{359}、2 回目は \ang{1}。
算術平均をとると
\[
  \frac{359 + 1}{2} = \ang{180}
\]
となり、\textbf{真南}という答えが返る。実際には 2 回とも北風だったのに。
\end{whybox}

正しいやり方は、角度を平均しないことである。
各時刻の風をまず東西成分 $u$ と南北成分 $v$ に分解し、
\textbf{成分を平均して}から、最後に一度だけ向きに戻す。

\begin{align}
  \bar{u} &= \frac{1}{N}\sum_i u_i, \qquad \bar{v} = \frac{1}{N}\sum_i v_i \\
  \theta &= \left(\arctan2(\bar{u}, \bar{v}) \cdot \frac{180}{\pi} + 180\right) \bmod 360
\end{align}

本パイプラインでは、この規則をコメントではなく\textbf{構造}で守っている。
統計を集計するモジュールは $u$ と $v$ しか扱わず、
風向という量は出力を書く直前の変換関数にしか現れない。
集計器に角度が入らないので、角度を平均するコードが書けない。

\paragraph{風向の向きの取り方}
本報告の風向はすべて\textbf{気象学の慣習}、すなわち
「風が吹いてくる方位」を北 \ang{0} から時計回りに測る。
西風（西から東へ吹く風、$u > 0$）は \ang{270} である。
既存シミュレータの \texttt{met\_to\_uv} と同じ定義で、両者は数値的に一致する。

\subsection{平均風速が 2 種類ある理由}
\label{sec:twomeans}

平均風速には意味の違う 2 つがあり、両方を出している。

\begin{description}
  \item[ベクトル平均風速 $|\bar{\mathbf{V}}|$]
    まず $u, v$ を平均し、その長さを取る：$\sqrt{\bar{u}^2 + \bar{v}^2}$。
    これは\textbf{正味どれだけ運ばれるか}である。
  \item[スカラー平均風速 $\overline{|\mathbf{V}|}$]
    各時刻の風速を先に計算してから平均する：$\frac{1}{N}\sum_i \sqrt{u_i^2 + v_i^2}$。
    これは\textbf{ふつうどれくらい強く吹いているか}である。
\end{description}

両者は一般に一致しない。数学的に必ず
\begin{equation}
  |\bar{\mathbf{V}}| \le \overline{|\mathbf{V}|}
\end{equation}
が成り立つ（三角不等式）。

\begin{whybox}{ビンの半分が東風、半分が西風だったら}
どちらも \SI{20}{\metre\per\second} だったとする。

スカラー平均は \SI{20}{\metre\per\second}。強い風が吹いていたのは事実である。

ベクトル平均は \SI{0}{\metre\per\second}。ビン全体をならして見れば、機体はどこにも運ばれない。

スペースバルーンの平均的な移流を見たいならベクトル平均、
そのビンで遭遇する風の強さを見たいならスカラー平均。
どちらか一方だけでは足りない。
\end{whybox}

この 2 つの比を\textbf{定常度}（constancy）と呼ぶ。

\begin{equation}
  C = \frac{|\bar{\mathbf{V}}|}{\overline{|\mathbf{V}|}} \in [0, 1]
\end{equation}

$C = 1$ なら常に同じ向きに吹いており、平均風向はそのまま信用してよい。
$C$ が 0 に近いと風向が定まっておらず、\textbf{平均風向には意味がない}。

このため本報告の図では、$C < 0.3$ の領域の風向を灰色でマスクしている。
定まらない風向を自信ありげに描くことは、
このパイプラインが犯しうる最も有害な誤りだからである。

\subsection{分散楕円と、それが実態を表さない場合}
\label{sec:ellipse}

風のばらつきは、$u$ と $v$ の分散共分散行列
\begin{equation}
  \Sigma = \begin{pmatrix} \sigma_u^2 & \sigma_{uv} \\ \sigma_{uv} & \sigma_v^2 \end{pmatrix}
\end{equation}
に全部入っている。この固有値分解から $1\sigma$ の楕円（長軸・短軸・傾き）が出る。
これはモンテカルロ計算で風に与える摂動の大きさを、推測ではなく実測から決められることを意味する。

NetCDF には楕円ではなく $\sigma_u^2, \sigma_v^2, \sigma_{uv}$ の 3 つを保存している。
楕円は情報を落とさずにここから復元できるが、逆はできないからである。

ただし楕円は\textbf{$(u,v)$ の散らばりが 2 変量正規に近い}ことを仮定している%
\footnote{%
「2 変量正規」は、2 つの量が同時にとる値の散らばり方の型のひとつである。
ここでの 2 つの量は東西成分 $u$ と南北成分 $v$ で、
10 年ぶんの解析を $(u,v)$ の点として平面に打つと散布図ができる。
その散布図が\textbf{中心が 1 つだけの、楕円形の雲}になっている状態を指す。
中心から離れるほど点がなだらかに減り、
どの向きに切っても釣鐘型（正規分布）の断面になる。

この型であれば、雲の形は
平均 2 つ（$\bar u$, $\bar v$）と分散共分散行列 $\Sigma$ の 3 つ、
合わせて 5 個の数だけで完全に決まる。
だからこそ楕円 1 つで散らばりを表せる。逆に言えば、
\textbf{楕円で表せるのは、この型に近い散らばりだけ}である。

破れ方は本文のとおりで、
中心が 2 つある（＝雲が 2 つの塊に分かれる）場合が本報告では実際に起きる。
2 つの塊の中間に中心を 1 つ置いた楕円は、
どちらの塊にも入らない場所を「典型」として描いてしまう。
中心が 1 つでも、裾が正規分布より重い（極端な値が想定より多い）場合には
楕円が \SI{90}{\percent} の範囲を過小に見積もる。
本報告が分位点（p90, p95）を別途保存しているのは、後者への備えでもある。
分位点は分布の型を仮定しないので、裾が重くてもそのまま正しい。%
}。
これが破れる場面がはっきりある。
日本上空の成層圏では、風向が 4--5 月と 9--10 月に東西逆転する。
その時期のビンには「強い西風の日」と「強い東風の日」が混ざり、
散らばりは 2 つの塊に分かれる。
このとき平均はほぼゼロ、分散は巨大、楕円は実際には誰も経験しない中間の風を描く。

この二峰性は、破裂高度帯で、しかも毎年決まった時期に起きる。
そのため風配図を用意した（図\ref{fig:f7-hokkaido-h1} 以降）。
風配図は分布そのものを見せるので、
2 つの塊があればそう見える。平均と楕円には見せられない。

\subsection{なぜ分位点も出すのか}

平均だけでは運用判断ができない。
「平均 \SI{15}{\metre\per\second}」でも、
そのビンの 1 割の日が \SI{40}{\metre\per\second} なら、回収計画は変わる。

そこで風速の経験分位点 $p_{10}, p_{25}, p_{50}, p_{75}, p_{90}, p_{95}$ と最大値、
さらに \SI{10}{}、\SI{20}{}、\SI{30}{}、\SI{40}{\metre\per\second} を超える頻度を保存した。

分位点は\textbf{風速（スカラー）に対してのみ}計算している。
風向には大小の順序がないので、分位点という概念が定義できない
（どこを起点に数えるかで答えが変わってしまう）。

なお分位点は平均や分散から計算することができない。
全標本を持っていないと出せない量である。
そのため集計では 1 ビン分の全標本をメモリに載せて厳密に計算している
（北海道の 1 ビンで約 \SI{46}{\mega\byte}、十分に載る）。
近似アルゴリズムは使っていない。

\subsection{600 個の観測は 600 個分の情報ではない}
\label{sec:neff}

H1 のビンには 10 年 $\times$ 15 日 $\times$ 4 回 $=$ 600 個の解析値が入る。
しかしこれは\textbf{独立な観測 600 個}ではない。

6 時間しか離れていない 2 つの解析値は、ほとんど同じ天気を見ている。
低気圧や高気圧が日本上空を通過するには 2--5 日かかるので、
15 日のビンに入る「独立な天気」はせいぜい 3--7 個である。

これを無視して標準誤差を $\sigma/\sqrt{600}$ とすると、
確からしさを 3 倍ほど過大に見積もることになる。

そこで各年の連続した区間内でラグ 1 自己相関 $r_1$ を求め、年をまたいで平均し、
\textbf{有効標本数}
\begin{equation}
  n_\mathrm{eff} = n\,\frac{1 - r_1}{1 + r_1}
\end{equation}
を保存した \citep{wilks2011statistical}。
実際の値は中央値で 76--78、$n$ の 13 \% である（3 領域とも）。
本報告の誤差はすべて $n_\mathrm{eff}$ から計算している。

\begin{whybox}{なぜ年をまたいで自己相関を取らないのか}
10 年分を時間順に並べて一括で自己相関を取ると、
1 月末の標本と翌年 1 月頭の標本が「隣」として扱われてしまう。
両者は 6 時間ではなく 1 年離れており、当然ほぼ無相関である。
その偽の無相関が混ざると $r_1$ が小さく出て、有効標本数を過大評価する。
だから各年の連続区間の中だけで計算し、あとで平均している。
\end{whybox}

\subsection{時刻の扱い}

JRA-3Q の解析時刻は 00, 06, 12, 18 UTC で、日本時間ではそれぞれ
09, 15, 21, 03 JST にあたる（18 UTC は JST では翌日になる）。

日付の区切りは既定で\textbf{UTC} を使った。
月をまたぐ繰り越しが不要で、データ本来の軸だからである。
7 日のビンでは、JST で区切った場合と食い違うのは 28 標本のうち 1 つだけである。

とはいえ主張で済ませるのは望ましくないので、両方を計算して比べた。
北海道で UTC 区切りと JST 区切りを突き合わせると、次のようになる。

\begin{center}
\begin{tabular}{lrr}
\toprule
量 & 最大の差 & RMS \\
\midrule
平均風速 & \SI{0.73}{\metre\per\second} & \SI{0.11}{\metre\per\second} \\
東西風 $\bar u$ & \SI{0.72}{\metre\per\second} & \SI{0.13}{\metre\per\second} \\
定常度 & 0.061 & 0.006 \\
飛行距離 & \SI{1.6}{\kilo\metre} & --- \\
標本数 & 10 個（400 個中） & --- \\
\bottomrule
\end{tabular}
\end{center}

平均風速で最大 \SI{0.73}{\metre\per\second}、
飛行距離では平均 \SI{130}{\kilo\metre} に対して最大 \SI{1.6}{\kilo\metre} である。
本報告のどの結論も、この選択では変わらない。

また、下層には本物の日変化がある（海陸風、夜間の接地逆転）。
そこで \SI{300}{\hecto\pascal} 以下の層に限り、
解析時刻ごとに分けた副次的な統計も出している。
図\ref{fig:f8-hokkaido} が示すとおり、夜（03 JST）と昼（15 JST）の平均風速の差は
3 領域・4 季節のどれでも \SI{1}{\metre\per\second} 未満で、日々の変動（数十 \si{\metre\per\second}）に比べて小さい。
それが\textbf{上層では時刻を区別せず平均してよい根拠}になっている。
これも主張ではなく、図で示している。

\subsection{領域平均の取り方}

緯度経度格子をそのまま平均すると、高緯度側の格子を過大に数えてしまう
（同じ \ang{0.375} でも、北の格子は地上での東西幅が狭い）。
JRA-3Q のファイルにはガウス求積の重みが入っているので、それを使って重み付き平均をとった。

\subsection{風から落下位置へ}
\label{sec:drift}

風の統計そのものより、多くの読者が知りたいのは「結局どこに落ちるのか」である。
そこで、既存シミュレータと同じ飛行モデルを風の統計に通し、
落下位置の気候値も作った。

飛行モデルは \texttt{backend/bfps/core/trajectory.py} と同一である。

\begin{itemize}
  \item 上昇は一定速度（既定 \SI{5}{\metre\per\second}）で破裂高度（既定 \SI{30}{\kilo\metre}）まで
  \item 降下は海面等価の降下率（既定 \SI{6}{\metre\per\second}）を空気密度でスケールした
        $v(z) = v_0\sqrt{\rho_0/\rho(z)}$
  \item 水平方向は慣性なしの純移流。機体はその高度の風にそのまま乗る
\end{itemize}

このモデルのもとでは、水平変位は
\begin{equation}
  \Delta x = \sum_k u_k T_k, \qquad \Delta y = \sum_k v_k T_k
  \label{eq:drift}
\end{equation}
と書ける。$T_k$ は気圧面 $k$ の担当する層に、上りと下りで滞在する合計時間である。
$T_k$ は層の厚さと降下則だけで決まり、\textbf{風には依存しない}。

\begin{whybox}{線形であることから何が言えるか}
式 \eqref{eq:drift} は風のプロファイルについて\textbf{線形}である。
そこから 2 つのことが従う。

第一に、\textbf{各解析時刻ごとに着地点を計算してから統計をとれる}。
落下位置の散らばりは推定ではなく実測になる。
もし平均風のプロファイルだけから散らばりを出そうとすると、
「高度 10 km の風のゆらぎと高度 25 km の風のゆらぎがどれくらい連動しているか」
を仮定しなければならない。その仮定は本文のどこにも現れないまま結果を左右してしまう。
層間の相関は保存していないので、仮定を置かずに各時刻の着地点を数えるほうを選んだ。

第二に、\textbf{平均風から求めた変位と、変位の平均が一致する}。
気候値から計算した量では一般には成り立たない性質で、
ここでは成り立つと言い切れるので、そう書いておく。
\end{whybox}

降下が寄与する時間は上昇より短い。
高度 \SI{30}{\kilo\metre} の空気は薄く、密度スケーリングによって降下速度が海面の 10 倍近くになるためである。
既定のパラメータでの飛行時間は約 \num{140} 分（上昇 100 分、降下 40 分）になる。

\subsubsection{\texorpdfstring{$T_k$}{T\_k} をどう計算するか}
\label{sec:drift-layers}

式\eqref{eq:drift} は $T_k$ が分かれば解ける。その $T_k$ の作り方を書いておく。

\paragraph{1. 気圧面に高度を与える}
JRA-3Q はジオポテンシャル高度を別の変数として持っているので、それを取得してある。
ジオポテンシャル高度は重力加速度の緯度・高度変化を吸収した量なので、
幾何学的な高度に直してから使う（\ref{sec:altitude} 節）。
高度は日々数百メートル動くが、層の厚さを決めるのが目的なので、
ビンごとの平均を 1 本使う。

\paragraph{2. 層の境界を決める}
気圧面 $k$ が「担当する」層は、\textbf{隣の気圧面との中点}で区切る。
高度の昇順に並べた $z_k$ に対し、

\begin{equation}
  z_k^{\text{下}} = \frac{z_{k-1} + z_k}{2}, \qquad
  z_k^{\text{上}} = \frac{z_k + z_{k+1}}{2}
  \label{eq:layer-edges}
\end{equation}

とし、これを打ち上げ高度と破裂高度で切り詰める。
層の厚さは $\Delta z_k = \max(z_k^{\text{上}} - z_k^{\text{下}},\, 0)$ で、
破裂高度より完全に上にある気圧面は $\Delta z_k = 0$ になって寄与しない。

\paragraph{3. 上りと下りの滞在時間を足す}
上昇は一定速度、降下はその層の代表高度での密度スケール速度を使う。

\begin{equation}
  T_k = \underbrace{\frac{\Delta z_k}{w_{\uparrow}}}_{\text{上り}}
      + \underbrace{\frac{\Delta z_k}{v_0\sqrt{\rho_0/\rho(z_k)}}}_{\text{下り}}
  \label{eq:layer-seconds}
\end{equation}

実装は \texttt{stats/drift.py} の \texttt{layer\_seconds} である。

\paragraph{実際の数字}
和歌山の中心格子、1 月前半の平均高度プロファイルで計算すると次のようになる。

\begin{table}[H]
\centering
\caption{$T_k$ の計算例（和歌山中心、1 月前半、既定の飛行諸元）。
$v_{\downarrow}$ はその高度での降下速度。}
\label{tab:layer-seconds}
\sisetup{table-format=5.0}
\begin{tabular}{@{}S[table-format=4.0]SSS[table-format=3.1]S[table-format=3.1]S[table-format=3.1]S[table-format=3.1]@{}}
\toprule
{気圧面} & {高度} & {層の厚さ} & {$v_{\downarrow}$} & {上り} & {下り} & {$T_k$} \\
{[\si{\hecto\pascal}]} & {[\si{\metre}]} & {[\si{\metre}]} & {[\si{\metre\per\second}]}
& {[\si{\second}]} & {[\si{\second}]} & {[\si{\second}]} \\
\midrule
1000 &   184 &  103 &  6.1 &  20.7 & 17.1 &  37.7 \\
 850 &  1491 &  234 &  6.5 &  46.7 & 36.2 &  82.9 \\
 500 &  5570 &  740 &  8.0 & 147.9 & 92.6 & 240.5 \\
 200 & 11894 &  821 & 11.8 & 164.2 & 69.5 & 233.6 \\
 100 & 16271 & 1188 & 16.7 & 237.7 & 71.3 & 309.0 \\
  30 & 23787 & 2234 & 30.4 & 446.8 & 73.6 & 520.4 \\
  10 & 31018 & 1281 & 53.5 & 256.3 & 23.9 & 280.2 \\
   5 & 35779 &    0 & 77.8 &   0.0 &  0.0 &   0.0 \\
   3 & 39417 &    0 &102.5 &   0.0 &  0.0 &   0.0 \\
\bottomrule
\end{tabular}
\end{table}

38 層ぶんを足すと $\sum_k T_k = \SI{8223}{\second}$（\num{137.0} 分）になり、
上昇が \num{99.4} 分、降下が \num{37.7} 分である。

表から読み取れることがいくつかある。

\begin{itemize}
  \item \textbf{上の層ほど厚い。} \SI{30}{\hecto\pascal} の層は \SI{2234}{\metre} あり、
        \SI{1000}{\hecto\pascal} の \SI{103}{\metre} の 20 倍を超える。
        気圧は高度に対して指数的に減るので、等間隔の気圧面は上ほど高度間隔が広くなる。
  \item \textbf{下りは上の層でほとんど時間を使わない。} \SI{30}{\hecto\pascal} では
        降下速度が \SI{30.4}{\metre\per\second} まで上がるので、
        \SI{2234}{\metre} を \SI{74}{\second} で通過する。
        同じ層を上るには \SI{447}{\second} かかる。
        成層圏の強い風が落下点を大きく動かすのは、主に\textbf{上りの側}である。
  \item \textbf{破裂より上の 3 層は寄与しない。} \SI{7}{\hecto\pascal}（約 \SI{33}{\kilo\metre}）、
        \SI{5}{\hecto\pascal}、\SI{3}{\hecto\pascal} は $T_k = 0$ になる。
        取得はしてあるが、既定の破裂高度 \SI{30}{\kilo\metre} では使われない。
        破裂高度を上げれば自動的に効きはじめる。
\end{itemize}

\begin{whybox}{最下層は地面ではなく \SI{1000}{\hecto\pascal} から始まる}
表の最下行を見ると、層の下端が \SI{184}{\metre} になっている。
最も低い気圧面そのものの高度であって、地面ではない。
\SI{1000}{\hecto\pascal} より下の空気は、この計算に入っていない。

飛行時間 \SI{8223}{\second} のうち \SI{40}{\second} 弱にあたる。
しかもそこは飛行中でいちばん風が弱く、いちばんゆっくり通る区間である。
落下点への影響は無視できる大きさなので、このままにしてある。

打ち上げ地点の標高が高い場合は、逆にこの層をまたいで始まることになる。
本報告の対象は海岸に近い低地なので問題にならないが、
山地から上げる計画に使うときは \texttt{launch\_altitude\_m} を設定する必要がある。
\end{whybox}

$T_k$ が風に依存しないことは、この 3 段階のどこにも風が出てこないことから確かめられる。
高度・層の厚さ・降下則だけで決まる。
そのため $T_k$ はビンごとに 1 回計算すればよく、
各解析時刻ではその重みに風を掛けて足すだけで着地点が出る。
季節による違いも小さく、和歌山では 1 月 \SI{8223}{\second} に対し 7 月 \SI{8259}{\second} で、
\SI{0.3}{\percent} しか変わらない。

このモデルには 2 つの限界があり、どちらも\textbf{既存シミュレータが置いている仮定と同じ}である。

\begin{itemize}
  \item 飛行中ずっと風のプロファイルが空間的にも時間的にも凍結していると見なしている。
        \texttt{wind.py} が打ち上げ点で 1 回取得したプロファイルを飛行中固定するのと同じ扱いである。
  \item 破裂高度を固定している。したがってここで出る散らばりは
        \textbf{風だけに由来する}もので、気球が早く割れたり粘ったりする効果は含まない。
\end{itemize}

\paragraph{向きの規約に注意}
落下位置の方位は「機体が\textbf{進んでいく}方位」で表す。
風向（風が\textbf{吹いてくる}方位）とはちょうど \ang{180} 逆である。
西風は \ang{270} だが、機体を運ぶ先は \ang{90} である。
両者を取り違えると着地点が打ち上げ地点の反対側に出てしまうので、
コードでも本報告でも別々の名前を使っている
（\texttt{dir\_vector\_mean} と \texttt{drift\_bearing\_deg}）。

\subsection{取得したデータが本当に正しいかの検証}

このパイプラインで最も重篤かつ最も気づきにくい失敗は、
\textbf{取り出した領域がずれていること}である。
緯度軸が北から南へ並んでいるため、範囲指定を逆向きに書くと、
例外は出ずに空の配列が返る。あるいは南半球の同じ形の領域が返ってくる。
どちらの場合も、その後の計算はすべて正常に完走し、
「もっともらしいが間違った」数字の集合ができあがる。

そのため、ダウンロードした全ファイル（2160 個）に対して、次を毎回検査している。

\begin{itemize}
  \item 座標が要求した範囲に格子 1 個分の余裕内で収まっているか
  \item 時刻数が期待通りで、6 時間刻みに欠落がないか
  \item 緯度が昇順に直っているか
  \item $|u|, |v| < \SI{200}{\metre\per\second}$、ジオポテンシャル高度が妥当な範囲か
  \item 全欠測でないか、欠測率が 1 \% 未満か
\end{itemize}

検査の結果はファイルごとの付随 JSON に記録してあるので、
後から「本当に全ファイルで検査が走ったか」を監査できる。

\begin{whybox}{検査で見つかった実例}
最初の試験取得で、1 月のジオポテンシャル高度に $-73$ gpm という負の値が出て検査に引っかかった。
調べると、これは\textbf{正しいデータ}だった。
発達した低気圧の中心では \SI{1000}{\hecto\pascal} 面が海面より下に潜る。
検査の下限を 0 にしていた側が間違っていたので、そちらを直した。

検査は間違いを見つけるためだけでなく、
自分の思い込みを見つけるためにも働く。
\end{whybox}
